%@ pages 94-107
%@ status translated
% Original section 10, "Critères pour qu'un préchamp soit un champ", pp. 94--107.
%
% Translation notes:
%  * The French abbreviates "n-préchamp" as "n-précat" in two places
%    (proof of Lemma \ref{grest} and proof of Proposition \ref{newcritere});
%    we read it as "n-prestack".
%  * In the paragraph defining \underline{Fl}^i(A;x,y) the original prints
%    "$s(x)=s(x)$ and $b(x)=b(y)$"; we reproduce the print (see the %
%    comment at that point).

\clearpage
\section{Criteria for a prestack to be a stack}
\label{caracterisationpage}

We begin by introducing several equivalent notions of {\em effectivity of
descent data}. Subsequently we shall use only version (ii) below; the others
are there to illuminate the definition---the reader may therefore skip the
proof.

The interest of conditions (vii) and (viii) below lies in their elementary
character: indeed, they refer neither to the model structure \ref{cmf}, nor to
the operation $SeCat$. However, without the cmf structure these conditions
would be practically unusable. We have included them to facilitate, if need
be, an exposition of the notion of stack without proofs.

\begin{lemmedefinition}
\label{lemmedef}
Let $A$ be a Segal $n$-prestack on a site $\Xx$. Let $X\in \Xx$ and let
$\Bb \subset \Xx /X$ be a sieve covering $X$. We shall say that {\em the
descent data for $A$ on $\Bb$ are effective} if one of the following
equivalent conditions is satisfied.
\newline
(i)\, Every morphism from $\ast _{\Bb}$ to $A|_{\Bb}$ in the homotopy
category $Ho(nSePCh(\Bb ^{\rm gro}))$ of Segal $n$-prestacks with respect to
object-by-object equivalences extends to a morphism from $\ast$ to $A(X)$ in
$Ho(nSePC)$. In other words, the natural map
$$
[\ast , A(X)]\rightarrow [\ast _{\Bb} , A|_{\Bb}]_{\Bb ^{\rm gro}}
$$
is surjective;
\newline
(ii)\, For every fibrant replacement $A\rightarrow A'$ for the coarse
topology, every morphism $\ast _{\Bb}\rightarrow A'$ extends to a morphism
$\ast _{\Xx /X}\rightarrow A'$;
\newline
(ii')\, There exists a fibrant replacement $A\rightarrow A'$ for the coarse
topology such that every morphism $\ast _{\Bb}\rightarrow A'$ extends to a
morphism $\ast _{\Xx /X}\rightarrow A'$;
\newline
(iii)\, For every fibrant replacement $A\rightarrow A'$ for the coarse
topology, the morphism of Segal $n$-categories
$$
A'(X)=\Gamma (\Xx /X , A'|_{\Xx /X})\rightarrow \Gamma (\Bb , A'|_{\Bb})
$$
is essentially surjective;
\newline
(iii')\, There exists a fibrant replacement $A\rightarrow A'$ for the coarse
topology such that the morphism of Segal $n$-categories
$$
A'(X)=\Gamma (\Xx /X , A'|_{\Xx /X})\rightarrow \Gamma (\Bb , A'|_{\Bb})
$$
is essentially surjective;
\newline
(iv)\, The morphism
$$
A(X)\rightarrow \lim _{\leftarrow , \Bb ^o} A|_{\Bb }
$$
is essentially surjective;
\newline
(v)\, For every fibrant replacement $A\rightarrow A'$ for the structure HBKQ
(\ref{cmfHirschho}), i.e. an object-by-object equivalence with each $A'(X)$
fibrant, and for every morphism
$$
\eta : D_{\Bb}\rightarrow A'
$$
there exists a homotopy $h: D_{\Bb}\times \overline{I}\rightarrow A'$ with
$h(0)=\eta$ and such that $h(1)$ comes from a morphism $\ast\rightarrow
A(X)$. Here $D_{\Bb}\rightarrow \ast _{\Bb}$ is the HBKQ-cofibrant
replacement constructed at the end of \S 5.
\newline
(v')\, There exists a fibrant replacement $A\rightarrow A'$ for the structure
HBKQ (\ref{cmfHirschho}) such that for every morphism
$$
\eta : D_{\Bb}\rightarrow A'
$$
there exists a homotopy $h: D_{\Bb}\times \overline{I}\rightarrow A'$ with
$h(0)=\eta$ and such that $h(1)$ comes from a morphism $\ast\rightarrow
A(X)$.

If we assume, moreover, that $A(Y)$ is a Segal $n$-category for every
$Y\in \Xx$, then these conditions are equivalent to the following conditions:
\newline
(vi)\, For every diagram of presheaves of Segal $n$-categories
$$
\ast _{\Bb } \leftarrow C \rightarrow A
$$
in which the first morphism is an object-by-object equivalence, there exists
$\ast \rightarrow A(X)$ such that the composite
$$
C\rightarrow \ast _{\Bb} \rightarrow \ast _{\Xx /X} \rightarrow A
$$
is homotopic (in Quillen's sense) to the original morphism $C\rightarrow A$;
\newline
(vii) \, Every diagram of presheaves of Segal $n$-categories
$$
\ast _{\Bb } \leftarrow C \rightarrow A
$$
in which the first morphism is an object-by-object equivalence can be
completed into a diagram
$$
\begin{array}{ccccc}
\ast _{\Bb }& \leftarrow &C &\rightarrow &A\\
\downarrow &&\downarrow && \parallel \\
\ast _{\Xx /X }& \leftarrow &C' &\rightarrow &A
\end{array}
$$
in which the morphism $C'\rightarrow \ast _{\Xx /X}$ is an object-by-object
equivalence; and
\newline
(viii)\, Let $nSeCat^{\Bb^o}$ be the category of presheaves of Segal
$n$-categories on $\Bb$; and let $Ho( nSeCat^{\Bb^o})$ be its Gabriel--Zisman
localisation, obtained by inverting the object-by-object equivalences. Then
every morphism $\ast _{\Bb}\rightarrow A$ in $Ho( nSeCat^{\Bb^o})$ comes from a
morphism $\ast \rightarrow A(X)$.
\end{lemmedefinition}
{\em Proof:}
\newline
(i)$\Leftrightarrow$(ii)($\Leftrightarrow$(ii'))\,\,
According to Quillen \cite{Quillen} every element of
$[\ast _{\Bb}, A|_{\Bb}]_{\Bb ^{\rm gro}}$ comes from a morphism
$f:\ast _{\Bb}\rightarrow A'|_{\Bb }$ (note that $A'|_{\Bb }$ is fibrant for
the coarse topology, see \S 4, and that $\ast_{\Bb}$ is automatically
cofibrant). If $f$ extends to $\ast _{\Xx /X}$ then the class of $f$ obviously
comes from $[\ast , A(X)]$. Conversely, if $f$ corresponds to an element of
$[\ast _{\Bb}, A|_{\Bb}]_{\Bb ^{\rm gro}}$ coming from $[\ast , A(X)]$ then
there exists a morphism $g:\ast _{\Xx /X}\rightarrow A'$ such that
$g|_{\Bb}$ is homotopic (on the left) to $f$ in Quillen's sense
\cite{Quillen}. Since $\ast _{\Bb} \times \overline{I}$ is a cylinder object
for $\ast_{\Bb}$, there exists a morphism $$
h: \ast _{\Bb} \times \overline{I}\rightarrow A'
$$
with $h(0)= f$ and $h(1)= g|_{\Bb}$. Now we have a morphism
$$
(\ast _{\Bb} \times \overline{I})\cup^{\ast _{\Bb} \times \{ 1\}}
(\ast _{\Xx /X} \times \{ 1\}) \rightarrow A'
$$
which extends along the trivial cofibration (for the coarse topology)
$$
(\ast _{\Bb} \times \overline{I})\cup^{\ast _{\Bb} \times \{ 1\}}
(\ast _{\Xx /X} \times \{ 1\}) \hookrightarrow
\ast _{\Xx /X} \times \overline{I},
$$
to a morphism whose restriction to $\ast _{\Xx /X} \times \{ 0\}$ provides the
desired extension of $f$.
\newline
(ii)$\Leftrightarrow$(iii)($\Leftrightarrow$(iii')) \, is immediate in view of
the definition of $\Gamma$.
\newline
(iii)$\Leftrightarrow$(iv) \, follows from Proposition \ref{calclim} below.
\newline
(i)$\Leftrightarrow$(v)($\Leftrightarrow$(v')) One argues as for
(i)$\Leftrightarrow$(ii) but with the HBKQ structure; cf. \S 5. Note that the
notion of homotopy class of maps is independent of the choice of cmf
structure and depends only on the weak equivalences; in particular in (i) the
sets in question are the same in the structure of \ref{cmf} and in the HBKQ
structure \ref{cmfHirschho}.

We now assume that $A$ is a presheaf of Segal $n$-categories. We have the
equivalence $Ho(nSeCat^{\Bb})\cong Ho(nSePCh(\Bb ^{\rm gro})_f)$ since there is
a natural ``fibrant replacement'' equivalence from $nSeCat^{\Bb}$ to
$nSePCh(\Bb ^{\rm gro})_f$. On the other hand we have the equivalence
$Ho(nSePCh(\Bb ^{\rm gro})_f) \cong Ho(nSePCh(\Bb ^{\rm gro}))$
(\cite{Quillen}). This gives (i)$\Leftrightarrow$(viii).
For conditions (vi) and (vii), we choose a fibrant replacement $A\rightarrow
A'$ for the coarse topology. Every morphism $\ast _{\Bb}\rightarrow A'$
therefore admits a factorisation
$$
\ast _{\Bb}\rightarrow E \rightarrow A'
$$
in which the first arrow is a trivial cofibration for the coarse topology and
the second a fibration for the coarse topology. There exists a (unique)
retraction $E\rightarrow \ast _{\Bb}$ (one may say that $\ast _{\Bb}$ is
fibrant for the coarse topology, but in fact it suffices to consider the
connected components of $C$). Set
$$
C:= E\times _{A'} A.
$$
The fact that the $A(Y)$ are Segal $n$-categories allows one to apply
\ref{proper} to obtain that $C\rightarrow E$ is an object-by-object
equivalence. Note, by \ref{stuff}, that the $C(Y)$ are Segal $n$-categories.
It follows that every element of $[\ast _{\Bb}, A|_{\Bb }]_{\Bb ^{\rm gro}}$
comes from a diagram (of presheaves of Segal $n$-categories)
$$
\ast _{\Bb}\leftarrow C
\rightarrow A
$$
in which the first morphism is an object-by-object equivalence.
We leave it to the reader to complete the proof that (vi) and (vii) are
equivalent to the other conditions, by expressing in the same way the
condition that such a diagram gives rise to an element of $[\ast _{\Bb},
A|_{\Bb }]_{\Bb ^{\rm gro}}$ which comes from $[\ast , A(X)]$.
\eop

\subnumero{Main criterion}

We have the following characterisation of stacks, which is analogous to the
definition of $1$-stack in Giraud \cite{GiraudThese} and
Laumon--Moret-Bailly \cite{LaumonMB}. Conditions (a) and (b) are the analogues
of the two conditions in the definition of a sheaf. Note also that condition
(a) leads us recursively to condition (b).


\begin{proposition}
\label{critere}
Let $n\geq 1$.
Let $A$ be a Segal $n$-prestack on $\Xx$ whose values are Segal $n$-categories.
Then $A$ is a Segal $n$-stack if and only if:
\newline
(a)\,\, for every $X\in \Xx$ and all $x,y\in A_0(X)$ the Segal
$(n-1)$-prestack $A_{1/}(x,y)$ on $\Xx /X$ is a Segal $(n-1)$-stack; and
\newline
(b) \, for every $X\in \Xx$ and every sieve $\Bb \subset \Xx /X$ in a set of
sieves which generates the topology, the descent data for $A$ with respect to
$\Bb$ are effective, see Definition \ref{lemmedef} above.
\end{proposition}
{\em Proof:}
Conditions (a) and (b) are obviously necessary (for (b) one uses the fact that
$A'$ is $\Gg$-fibrant by virtue of Lemma \ref{grofibchampimplfib}).
Suppose that (a) and (b) are satisfied and let $A'\rightarrow A''$ be a
$\Gg$-trivial cofibration to a $\Gg$-fibrant Segal $n$-prestack. For every
$X\in \Xx$ and $x,y\in A_0(X)$, the morphism
$$
A_{1/} (x,y)\rightarrow         A''_{1/}(x,y)
$$
is a $\Gg_X$-trivial cofibration of Segal $(n-1)$-prestacks on $\Xx /X$ whose
target is $\Gg _X$-fibrant (Corollary \ref{Gfibrestrict}). Condition (a)
implies that this morphism is an equivalence for the coarse topology, in
particular
$$
A(X)_{1/} (x,y)\rightarrow      A''(X)_{1/}(x,y)
$$
is an equivalence. The same holds for the morphism
$$
A'(X)_{1/} (x',y')\rightarrow  A''(X)_{1/}(x',y')
$$
for all $x',y'\in A'_0(X)$ since $x'$ and $y'$ are equivalent to objects
coming from $A_0(X)$. This shows that $A'\rightarrow A''$ is fully faithful
object-by-object. It remains to see that $A'\rightarrow A''$ is essentially
surjective object-by-object. For this, one may choose a factorisation
$$
A'\rightarrow B \rightarrow A''
$$
in which the first morphism is a trivial cofibration for the coarse topology
and the second a fibration for the coarse topology. Replacing $A'$ by $B$ if
necessary, we may henceforth assume that $A'\rightarrow A''$ is a fibration
for the coarse topology.

Let $X\in \Xx$ and $u\in A''_0(X)$, which we regard as a morphism
$$
u: \ast _X \rightarrow A'' |_{\Xx /X}.
$$
Set
$$
C:= (A'|_{\Xx /X})\times _{A'' |_{\Xx /X}}\ast _X .
$$
The morphism $C\rightarrow \ast _X$ is a fibration for the coarse topology and
also an equivalence for the topology $\Gg$. Moreover, since
$A'\rightarrow A''$ is fully faithful, the morphism $C\rightarrow \ast _X$ is
fully faithful. In particular, $C$ is $0$-truncated, and equivalent
object-by-object to a presheaf of subsets of $\ast _X$. Since our morphism is
an equivalence for the topology $\Gg$, there exists a sieve
$\Bb \subset \Xx /X$ of the topology (and even of the generating set of sieves
mentioned in (b)) such that $\ast _{\Bb}$ is contained in this subpresheaf of
$\ast _X$. In particular, the morphism
$$
C\times _{\ast _X}\ast _{\Bb} \rightarrow \ast _{\Bb}
$$
is an object-by-object equivalence and, since it is also a fibration for the
coarse topology, there exists a section $\ast _{\Bb }\rightarrow C$. We thus
obtain a morphism
$$
\ast _{\Bb} \rightarrow A'|_{\Xx /X}
$$
which, thanks to (b), extends to a morphism
$$
u':\ast _{X} \rightarrow A'|_{\Xx /X}.
$$
We thus obtain an element $u'\in A'_0(X)$. By the same type of argument one
sees that the image of $u'$ in $A''_0(X)$ is equivalent to $u$, which proves
essential surjectivity.
\eop

\subnumero{Criterion for Segal \texorpdfstring{$0$}{0}-stacks}

We have a result analogous to \ref{critere} for Segal $0$-prestacks, i.e.
simplicial presheaves. To obtain it, we use the following lemma.

\begin{lemma}
\label{pimseg}
Let $m,n\geq 0$ and let $A$ be a presheaf of Segal $n$-categories.
Then $A$ is a Segal $n$-stack if and only if $\Pi _{m,Se}\circ A$ is a Segal
$(n+m)$-stack.
\end{lemma}
{\em Proof:}
First we show that if $A$ is a Segal $n$-stack, then $\Pi _{m,Se}\circ A$ is a
Segal $(n+m)$-stack. Reasoning by induction on $m$ and applying criterion
\ref{critere}, we reduce to the case $n=0$ and $m=1$. So let $A$ be a Segal
$0$-stack. We may assume that $A$ is $\Gg$-fibrant. We shall prove conditions
(a) and (b) of \ref{critere} for $B=\Pi _{1,Se}\circ A$. Note that $A$ and
$\Re _{\geq 0}B$ are equivalent (object-by-object). If $X$ is an object of
$\Xx$ and $\Bb$ a sieve covering $X$, every morphism $\ast _{\Bb}\rightarrow B$
induces a morphism $\ast _{\Bb}\rightarrow \Re _{\geq 0}B$ and hence descent
data for $A$ (i.e. an element of $[\ast _{\Bb}, A|_{\Bb }]_{\Bb ^{\rm gro}}$).
Now the inclusion $\ast _{\Bb}\rightarrow \ast _{\Xx /X}$ is a
$\Gg$-equivalence and the fact that $A$ is $\Gg$-fibrant therefore implies the
existence of an extension of our descent data to an element of
$[\ast , A(X)]= [\ast , B(X)]$. This proves condition (b) of \ref{critere} for
$B$. On the other hand, let $a$ and $b$ be in $B_0(X)$. We have the equivalence
$$
B_{1/}(a,b) \cong Path ^{a,b}(A):= \ast \times _{A|_{\Xx /X}}\ast ,
$$
and this latter prestack is a Segal $0$-stack by \ref{fiprodchamp}. This gives
condition (a) of \ref{critere}, and hence, by \ref{critere}, $B$ is a Segal
$1$-stack.

For the other direction, suppose that $\Pi _{m,Se}\circ A$ is a Segal
$(n+m)$-stack. Choose a $\Gg$-fibrant replacement $A\rightarrow A'$. We obtain
a morphism
$$
\Pi _{m,Se}\circ A\rightarrow \Pi _{m,Se}\circ A',
$$
whose target is a Segal $(n+m)$-stack by the first paragraph. By hypothesis the
source is a Segal $(n+m)$-stack. Since this morphism is a weak
$\Gg$-equivalence, \ref{morphentrechamps} implies that it is an
object-by-object equivalence. This in turn implies that $A\rightarrow A'$ is an
object-by-object equivalence, and hence $A$ is a Segal $n$-stack.
\eop


We now give a version of the criterion adapted to Segal $0$-prestacks (or
simplicial presheaves). Let us recall here that a simplicial presheaf $U$ is a
Segal $0$-stack if and only if $U$ is {\em flasque with respect to the objects
of the site} in Jardine's sense \cite{Jardine}, or {\em of cohomological
descent} in Thomason's terminology \cite{Thomason} \cite{TravauxThomason}.

\begin{corollary}
\label{0critere}
If $U$ is a simplicial presheaf over $\Xx$, then $U$ is a Segal $0$-stack if
and only if
\newline
(a)\,\, for every $X\in \Xx$ and all $x,y\in U_0(X)$ the simplicial presheaf
\newline
$Y\mapsto Path ^{x|_Y,y|_Y}(U(Y))$ on $\Xx /X$ is a Segal $0$-stack; and
\newline
(b)\,\, for every $X\in \Xx$ and every sieve $\Bb \subset \Xx /X$ (in a family
of sieves which generates the topology), the descent data for $U$ on $\Bb$ are
effective.
\end{corollary}
{\em Proof:}
One observes, as in the preceding proof, that if $A:= \Pi _{1,Se}\circ U$ then
one has
$$
\left( Y\mapsto Path ^{x|_Y,y|_Y}(U(Y))\right) \cong A_{1/}(x,y)
$$
and on the other hand, the set of descent data (resp. effective descent data)
for $U$ is isomorphic to the set of descent data (resp. effective descent
data) for $A$. One concludes with the help of \ref{critere} and \ref{pimseg}.
\eop


In condition (a) of this corollary, $Path ^{x|_Y,y|_Y}(U(Y))$ denotes the
simplicial set of paths between $x|_Y$ and $y|_Y$ in a Kan replacement of
$U(Y)$.


If $U$ is a $k$-truncated Segal $0$-stack (i.e. the $U(X)$ have vanishing
homotopy groups in degree $>k$) then the criterion of \ref{0critere} becomes
recursive, since in condition $(a)'$, $Path ^{x|_Y,y|_Y}(U(Y))$ is
$k-1$-truncated. One can convert this remark into a version of Proposition
\ref{critere} adapted to $k$-stacks (plain, and not Segal).

\medskip

\subnumero{Compatibility with restriction to \texorpdfstring{$\Xx /X$}{X/X}}

We return to Segal $n$-stacks; we can now complete the results of \S 4
concerning the topology $\Gg$.



\begin{lemma}
\label{grest}
Let $\Xx$ be a site whose Grothendieck topology we denote by $\Gg$. A Segal
$n$-prestack $A$ on $\Xx$ is a stack if and only if $A|_{\Xx /X}$ is a stack
for every $X\in \Xx$. The restriction of Segal $n$-prestacks
$$
p^{\ast}: A\mapsto A|_{\Xx /X}
$$
preserves $\Gg$-fibrations and trivial $\Gg$-fibrations; for the topology
$\Gg$, $p^{\ast}$ is a Quillen functor both on the left and on the right.
\end{lemma}
{\em Proof:}
It is obvious from our criteria (\ref{critere} and \ref{0critere}) that $A$ is
a stack if and only if $A|_{\Xx /X}$ is a stack for every $X\in \Xx$. In
particular the restriction $p^{\ast}$ sends $\Gg$-fibrant objects to
$\Gg$-fibrant objects (cf. Lemma \ref{grofibchampimplfib}). Let
$f: A\rightarrow B$ be a $\Gg$-fibration between $\Gg$-fibrant objects. Then
$f|_{\Xx /X}$ is a fibration for the coarse topology (by \ref{shriek}) between
$\Gg$-fibrant objects. By Lemma \ref{stuff}, $f|_{\Xx /X}$ is $\Gg$-fibrant.

Now we know that $p^{\ast}$ turns fibrations between fibrant objects into
fibrations. But a cofibration in a cmc is trivial if and only if it satisfies
the lifting property with respect to all fibrations between fibrant objects.
This implies that the adjoint $p_!$ of $p^{\ast}$ turns trivial cofibrations
into trivial cofibrations (and we already knew that $p_!$ preserves
cofibrations). Hence $p_!$ is a left Quillen functor, and consequently
$p^{\ast}$ is a right Quillen functor.
\eop



\subnumero{Concerning \texorpdfstring{$\emptyset$}{the empty set}}
% The authors abbreviate "n-précatégorie de Segal" as "n-précat de Segal";
% "précat" is not a slip for "préchamp".  "n-précat" is translated
% "Segal n-precategory", "n-préchamp" as "Segal n-prestack".
Usually in a site $\Xx$ one has an initial object $\iota$ (generally the
``empty'' object $\iota = \emptyset$), and the stacks on $\Xx$ must satisfy the
condition $A(\iota )\cong \ast$, at least when the topology $\Gg$ admits the
empty sieve $\Bb = \emptyset \subset \Xx /\iota$ as a sieve covering $\iota$
(which will be the case in the geometric examples). Indeed, under this
hypothesis, if $A$ is $\Gg$-fibrant, then since
$E'\times \ast _{\iota} \rightarrow E\times \ast _{\iota}$ is a trivial
cofibration, for every cofibration of Segal $n$-precategories
$E'\subset E$, every morphism $E'\rightarrow A(\iota )$ extends to
$E\rightarrow A(\iota )$, which entails that $A(\iota )\rightarrow \ast$ is a
trivial fibration.


Let us see what this gives with our criteria: since
$\Bb = \emptyset \subset \Xx /\iota$ is a sieve covering $\iota$ ($\Bb$ does
not even contain $\iota$), the prestack $\ast _{\Bb}$ is equal to the empty
presheaf $\emptyset$ over $\Xx /\iota$, so there exists a unique morphism
$\ast _{\Bb}\rightarrow A'|_{\Xx /\iota}$. Condition (b) of our criteria
therefore implies that $A'(\iota )$ is non-empty. On the other hand, condition
(a) of our criteria implies (by induction on $n$) that for each pair of
objects of $A(\iota )$ the corresponding $(n-1)$-category of morphisms is
$\ast$. Hence $A(\iota )=\ast$.
For $n=0$ the same argument (replacing $A$ each time by its
$\Pi _{1,Se}\circ A$) shows by induction that all the homotopy groups of the
simplicial set $A(\iota )$ vanish; this is the initial case of the induction
for $n>0$.
On the other hand, for what is called the ``coarse topology'', one does not
admit the empty sieve $\emptyset \rightarrow \Xx /\iota$ as a covering sieve.
In that case one does not have the condition $A(\iota )\cong \ast$.


\subnumero{Dévissage}



For $n>0$ one can break up the condition for being a stack in terms of the
$k<n$, thanks to the following corollary. Let us recall once again that a
simplicial presheaf $U$ is a Segal $0$-stack if and only if $U$ is flasque with
respect to the objects of the site (Jardine \cite{Jardine}), or of
cohomological descent (Thomason \cite{Thomason} \cite{TravauxThomason}). This
condition is therefore well known.

We recall (cf. \S 2) that if $A$ is a Segal $n$-category, for $k\leq n$ one
has at one's disposal the {\em $k$-groupic interior} $A^{int, k}$, which is a
Segal $k$-category. Essentially this construction consists in keeping only the
$i$-arrows which are invertible up to equivalence ($i>k$) and in regarding the
result as a Segal $k$-category. For $k=0$, $A^{int, 0}$ is therefore a
simplicial set with $\Pi _{n,Se}\circ A^{int, 0}$ equivalent to the largest
Segal $n$-groupoid contained in $A$. These constructions extend in the
obvious way to presheaves of Segal $n$-categories.

The following corollary allows one to break up the condition for a Segal
$n$-prestack $A$ to be a Segal $n$-stack in terms of the condition for a
simplicial presheaf to be a Segal $0$-stack.



\begin{corollary}
\label{devissage}
Let $A$ be a Segal $n$-prestack ($n\geq 1$) such that $A(X)$ is a Segal
$n$-category for every $X\in \Xx$. Then $A$ is a Segal $n$-stack if and only
if:
\newline
(a) \, for every $X\in \Xx$ and every pair of objects $x,y\in A_0(X)$, the
Segal $(n-1)$-prestack $A_{1/}(x,y)$ is a Segal $(n-1)$-stack over
$\Xx /X$; and
\newline
(b) \, $A^{int, 0}$ is a Segal $0$-stack.
\end{corollary}
{\em Proof:}
One has
$$
[\ast _{\Bb}, A|_{\Bb}]_{\Bb ^{\rm gro}} =
[\ast _{\Bb}, A^{int, 0}|_{\Bb}]_{\Bb ^{\rm gro}} ,
$$
and one applies \ref{critere} and \ref{0critere}.
\eop


\begin{corollary}
If $A$ is a Segal $n$-stack then for $0\leq k < n$, $A^{int,k}$ is a Segal
$k$-stack.
\end{corollary}
\eop


In fact, the reader may convince themselves that the present dévissage reduces
us to an infinite number of effectivity conditions for descent data, one for
the $i$-arrows for each $i\geq 0$. To state this more precisely, let us
introduce some notation. Let us first recall a notation of Tamsamani
\cite{Tamsamani}: the set of $i$-arrows of an $n$-category $A$, denoted
$Fl^i(A)$, is the set $A_{1^i}$ where $1^i = (1,\ldots , 1)\in \Theta ^n$ is
the class of $(1,\ldots , 1,0,\ldots , 0)\in \Delta ^n$. One has the maps $s$
(``source'') and $b$ (``target'') from $Fl^i(A)$ to $Fl^{i-1}(A)$. In our
situation, if $A$ is a Segal $n$-category, we shall agree to call the
$i$-arrows for $i>n$ the $i$-arrows of $\Pi _{m,Se}\circ A$ for $m> i-n$. More
systematically one could define (in view of our definition of $\tau _{\leq m}$,
see \S 2)
$$
Fl^i(A):= Fl^i(\tau _{\leq i+1}A).
$$
Let $i\geq 1$ and let $x,y\in Fl^{i-1}(A)$ with $s(x)=s(x)$ and $b(x)=b(y)$.
% the original prints "$s(x)=s(x)$" here; the same condition is printed
% correctly, as $s(x)=s(y)$, at the two later occurrences below. We reproduce
% the print.
One then defines the Segal $(n-i)$-category
$\underline{Fl}^i(A; x,y)$ as the sub-Segal-$(n-i)$-category of $A_{1^i/}$
consisting of the objects whose source is $x$ and whose target is $y$ (here, if
$n-i$ is negative, one replaces it by $0$).
\footnote{
This inelegance should eventually disappear when one can speak of the
{\em $\infty$-categories}.}

If $A$ is a Segal $n$-prestack whose values $A(X)$ are Segal $n$-categories,
then for every $X\in \Xx$ and $x,y\in Fl^{i-1}(A(X))$ with $s(x)=s(y)$ and
$b(x)=b(y)$, one has at one's disposal the Segal $(n-i)$-prestack
$\underline{Fl}^i(A|_{\Xx /X}; x,y)$ which to $Y\rightarrow X$ associates the
$(n-i)$-category $\underline{Fl}^i(A(Y); x|_Y,y|_Y)$.

\begin{proposition}
\label{conditiond}
Let $A$ be a Segal $n$-prestack on $\Xx$ such that the $A(X)$ are Segal
$n$-categories. Then $A$ is a Segal $n$-stack if and only if for every
$X\in \Xx$ and every sieve $\Bb \subset \Xx /X$ (in a family of sieves which
generates the topology), one has:
\newline
---the descent data for $A$ with respect to $\Bb$ are effective; and
\newline
---for every $1\leq i < \infty$, for all $x,y\in Fl^{i-1}(A(X))$ with
$s(x)=s(y)$ and $b(x)=b(y)$, the descent data for
$\underline{Fl}^i(A|_{\Xx /X}, x,y)$ on $\Bb$ are effective.
\end{proposition}
The proof is left as an exercise to the reader. {\em Hint:} combine Lemma
\ref{pourloc}, with the remark that it suffices, in order
to obtain all trivial cofibrations, to consider adequate sequences composed of
trivial cofibrations for the coarse topology, as well as cofibrations of the
form $cof(U\hookrightarrow U'; \Bb \subset \Xx /X)$ where $U'= h(1^i)$
represents the $i$-arrows, and $U$ is its ``boundary'' which represents the
pairs of $i-1$-arrows with the same source and target.
\eop



\subnumero{The case of Segal \texorpdfstring{$1$}{1}-prestacks and presheaves of
simplicial categories}

In the last part of the paper we shall be mainly interested in the case $n=1$;
for this reason we copy out \ref{devissage} for Segal $1$-categories as well as
for presheaves of simplicial categories. Let us recall once again that the
Segal $0$-stacks are exactly the simplicial presheaves flasque with respect to
the objects of the site \cite{Jardine}, or equivalently satisfying the
cohomological descent condition of \cite{Thomason} \cite{TravauxThomason}.

\begin{corollary}
Let $A$ be a Segal $1$-prestack such that $A(X)$ is a Segal category for every
$X\in \Xx$. Then $A$ is a Segal $1$-stack ($\infty$-stack) if and only if:
\newline
(a) \, for every $X\in \Xx$ and every pair of objects $x,y\in A_0(X)$, the
simplicial presheaf $A_{1/}(x,y)$ over $\Xx /X$ is a Segal $0$-stack; and
\newline
(b) \, the simplicial presheaf $A^{int, 0}$ is a Segal $0$-stack.
\end{corollary}
\eop

If $B$ is a simplicial category, one can compute the simplicial set
$B^{int, 0}$ in the following way: let $WB$ be the simplicial subcategory in
which one keeps only the arrows which are invertible up to homotopy; let
$\nu WB$ be its ``nerve'', i.e. the associated bisimplicial set (it is also
the associated Segal $1$-precategory). Then, if $d(\nu WB)$ is the corresponding
diagonal simplicial set, one has $d(\nu WB)\cong B^{int, 0}$.
This discussion extends in the obvious way to presheaves of simplicial
categories (which give rise to Segal $1$-prestacks denoted by the same
letter).

\begin{corollary}
Let $B$ be a presheaf of simplicial categories. Then $B$ is a Segal $1$-stack
if and only if:
\newline
(a) \, for every $X\in \Xx$ and every pair of objects $x,y\in Ob (B)(X)$, the
simplicial presheaf $Mor_B(x,y)$ over $\Xx /X$ is a Segal $0$-stack, and
\newline (b) \, the simplicial presheaf
$d(\nu WB)= B^{int,0}$ is a Segal $0$-stack.
\end{corollary}
\eop





\subnumero{A variant of the criterion}

%%%%%% this subsection modified in v3 including the comments at the end %%%%

We give here a statement very slightly different from \ref{critere}.
We remark first that if $A'$ is a Segal $n$-prestack fibrant for the coarse
topology, then $A'|_{\Xx /X}$ is also fibrant for the coarse topology, as is
$A'|_{\Bb}$ for every sieve $\Bb$ covering $\Xx /X$ (see \S 4). Moreover, we
observe that the morphism
$$
\Gamma (\Xx /X , A'|_{\Xx /X}) \rightarrow A'(X)
$$
is an isomorphism of Segal $n$-categories, since $X$ is the final object of the
category $\Xx /X$ and $\Gamma$ is simply the functor of global sections for
presheaves.

The following statement is probably not true for Segal $n$-prestacks, even for
$n=0$; this is due to the fact that Lemma \ref{pourloc} (corrected in {\tt v3})
is true only for the $n$-truncated case. It would seem that one obtains a
counterexample to \ref{newcritere} by considering a topological space $X$
having two non-empty open subsets $U$ and $V$ and no others apart from the
open subsets of $U\cap V$ introduced by the recurrence condition
$U\cap V \cong X$.

%%%% in v3 the following only applies to the truncated case %%%
%%%%
\begin{proposition}
\label{newcritere}
Let $A$ be an $n$-prestack (not Segal) on $\Xx$ whose values are $n$-categories.
Fix a trivial cofibration for the coarse topology $A\rightarrow A'$ with $A'$
fibrant for the coarse topology.
Then $A$ is an $n$-stack if and only if:
\newline
(c)\,\,
for every $X\in \Xx$ and every sieve $\Bb \subset \Xx /X$ of the topology,
the morphism
$$
A'(X)=\Gamma (\Xx /X, A'|_{\Xx /X}) \rightarrow \Gamma (\Bb , A'|_{\Bb})
$$
is an equivalence of $n$-categories.
\end{proposition}
{\em Proof:}
Suppose that $A$ is an $n$-stack. Replacing $A$ by $A'$ if necessary, we may
also assume that $A$ is fibrant for the coarse topology, and hence even
$\Gg$-fibrant (see \ref{stuff}). We shall show that, in this case, the morphism
of (c) is a trivial fibration; for this we shall show that it has the lifting
property with respect to every cofibration $E\rightarrow E'$ of $n$-precategories.

We begin by observing that if a morphism $U\rightarrow V$ of $n$-prestacks over
$\Xx /X$ is an equivalence over each object of a sieve $\Bb \subset \Xx /X$
covering $X$, then it is a weak $\Gg$-equivalence. Indeed, this is easy to see
for $n=0$ and for $n>0$ one may proceed by induction on $n$ using directly the
definition of weak $\Gg$-equivalence (\S 3). We apply this to the morphism
$$
j:E'_{\Bb} \cup ^{E_{\Bb}} E_{\Xx /X} \rightarrow E'_{\Xx /X},
$$
where $E\hookrightarrow E'$ is the given trivial cofibration of Segal
$n$-precategories, and $E_{\Bb}$ is the $n$-prestack obtained by applying $p_{!}$
to the constant $n$-prestack with values $E$ on $\Bb$, $p$ denoting the
inclusion $\Bb \rightarrow \Xx /X$. % the original abbreviates "n-préchamp"
% as "n-précat" here
We thus obtain that $j$ is a $\Gg$-trivial cofibration. The lifting property of
$A\rightarrow \ast$ for the morphism $j$ translates into the lifting property
of the morphism
$$
\Gamma (\Xx /X, A|_{\Xx /X}) \rightarrow \Gamma (\Bb , A|_{\Bb})
$$
with respect to the cofibration $E\rightarrow E'$. This proves that the
morphism in question is a trivial fibration, which gives condition (c) for $A$.

For the other direction, suppose that $A$ satisfies condition (c).
Note that if $A$ is fibrant for the coarse topology, condition (c) implies that
$A$ has the lifting property with respect to every cofibration of the form
$cof(U\hookrightarrow U', \Bb \subset \Xx /X)$ (see the definition just before
Proposition \ref{naqualocaliser}). By Lemma \ref{pourloc}, this implies that
$A$ has the lifting property with respect to every $\Gg$-trivial cofibration,
i.e. that $A$ is $\Gg$-fibrant. Condition (c) therefore implies that $A$ is an
$n$-stack.
\eop


%%%%%%%%%%%%%%%%%%
{\em Remarks:} We indicate here an alternative argument of induction on $n$
which allows one easily to reduce the proposition to the case $n=0$. Suppose
that $n\geq 1$ and that one has already proved the proposition for Segal
$(n-1)$-prestacks. If $\Bb$ is a sieve covering $\Xx /X$, then for all
$x,y\in A'_0(X)$ one has
$$
\Gamma (\Bb  , A'_{1/}(x,y)|_{\Bb})= \Gamma (\Bb , A'|_{\Bb})_{1/}(x,y).
$$
In view of the definition of the equivalence between Segal $n$-categories,
condition (c) for $A'$ is manifestly equivalent to the conjunction of condition
(b) of \ref{critere} with condition (c) for the $A'_{1/}(x,y)$. By the
induction hypothesis, the latter is equivalent to condition (a) of
\ref{critere}. One thus obtains indeed that condition (c) is equivalent to
conditions (a) plus (b) of \ref{critere}.
For $n=0$, i.e. for presheaves of sets, there is an easy argument.


In the case of simplicial presheaves, the fact that condition (c) is necessary
was the subject of a note by Toën \cite{ToenDescente}.
For every $n$ the condition is necessary for Segal $n$-prestacks.

It is probable that, replacing Lemma \ref{pourloc} by the version in terms of
hypercoverings proved recently by Dugger \cite{Dugger},
one would obtain a necessary and sufficient version of this criterion (c) in
terms of hypercoverings. We leave this interesting question open.




\subnumero{The proto-stacks}

For ordinary (non-Segal) $1$-stacks, there appears in Giraud \cite{GiraudThese}
and Laumon--Moret-Bailly \cite{LaumonMB} a notion of ``prestack'': in their
sense, a prestack is one of our prestacks which moreover satisfies condition
(a) of \ref{critere}. In the case of presheaves, this corresponds to requiring
the first of the two conditions of the usual definition of sheaves (Houzel
called these semi-sheaves, but ``semi-stack'' does not sound good) and the
terminology of \cite{GiraudThese} \cite{LaumonMB} is therefore not compatible
with the usual terminology concerning presheaves; this is why we have
preferred to keep the term ``prestack'' for objects devoid of any relation with
the topology $\Gg$. However, the notion of a prestack satisfying condition (a)
of \ref{critere} is useful, and we adopt the following definition:

\begin{definition}
Let $n\geq 1$.
A Segal $n$-prestack $A$ will be called a {\em protostack} if the $A(X)$ are
Segal $n$-categories, and if $A$ satisfies condition (a) of \ref{critere},
namely that for every $X\in \Xx$ and all $x,y\in A(X)_0$, the Segal
$(n-1)$-prestack $A_{1/}(x,y)$ is a Segal $(n-1)$-stack on $\Xx /X$.
\end{definition}


\begin{lemma}
Let $A$ be a Segal $n$-prestack whose values $A(X)$ are Segal
$n$-ca\-te\-go\-ries. Let $f:A\rightarrow A'$ be an associated stack
(\S 9). Then $A$ is a protostack if and only if the morphism $f$ is fully
faithful (object-by-object).
\end{lemma}
{\em Proof:}
If $f$ is fully faithful then, since the $A'_{1/}(fx,fy)$ are Segal
$(n-1)$-stacks, the same holds for the $A_{1/}(x,y)$. Conversely, suppose that
$A$ is a protostack, and let $x,y\in A(X)_0$.
Then the morphism $A_{1/}(x,y)\rightarrow A'_{1/}(fx,fy)$ is a
$\Gg_X$-weak equivalence between Segal $(n-1)$-stacks on $\Xx /X$, hence by
\ref{morphentrechamps} it is an object-by-object equivalence.
Hence $f$ is fully faithful.
\eop

The next proposition is the analogue of a result of \cite{flexible}. We leave
the proof as an ``exercise'' to the reader.

\begin{proposition}
\label{exo}
Let $F$ be a Segal $n$-protostack on $\Xx$, and define a Segal $n$-prestack $G$
by
$$
G(X):= {\rm colim} _{\Bb \subset \Xx /X}\left( \lim _{\leftarrow , \Bb ^o}
F|_{\Bb} \right)
$$
(here one uses the limits and colimits of \cite{limits}, see \S 14 for how to
compute them; the colimit is taken over the filtered set of sieves covering
$X$). Then $F\rightarrow G$ is the associated stack.
\end{proposition}
\eop

Recall that one may regard an $n$-category as a Segal $n$-category by
regarding the sets of $n$-morphisms as $0$-truncated simplicial sets. The Segal
$n$-categories which correspond in this way to $n$-categories are exactly those
which are {\em $n$-truncated}, i.e. for which one has
$A\cong A^{int, n}\cong \tau _{\leq n}A$. The following lemma indicates that one
may use this correspondence to translate the results on Segal $n$-prestacks
into results on (non-Segal) $n$-prestacks.
In this lemma, the hypothesis that the $A(X)$ must be Segal $n$-categories is
indeed necessary (since the operation $SeCat$ does not preserve the property of
being truncated).


\begin{corollary}
\label{preservetronque}
If $A$ is a presheaf of Segal $n$-categories which is $n$-truncated (i.e.
$A\cong \tau _{\leq n}A$) and if $A\rightarrow A'$ is an associated stack, then
$A'$ is also $n$-truncated.
\end{corollary}
{\em Proof:} This would be an immediate consequence of
Proposition \ref{exo}.
Since we have not provided a proof of that proposition, we indicate here an
alternative proof of the corollary.
One may suppose that, in parallel with our discussion of Segal $n$-prestacks,
one has developed an analogous discussion for non-Segal $n$-prestacks.
For example, for a non-Segal $n$-prestack $T$, one can define the associated
$n$-stack $f:T\rightarrow T'$. This morphism is a weak $\Gg$-equivalence of
$n$-prestacks. Let $Ind(T)$ (resp. $Ind(T')$) denote the induced Segal
$n$-prestacks.
Suppose that the $T(X)$ are $n$-categories.
The fact that $f$ is a $\Gg$-equivalence of $n$-prestacks implies that
$Ind(f)$ is a $\Gg$-equivalence of Segal $n$-prestacks. On the other hand, by
\ref{critere} and its analogue for non-Segal $n$-prestacks, $Ind(A')$ is a
Segal $n$-stack. Indeed, with condition (a) one reduces by induction to the
case $n=0$; for that case, the question is whether a sheaf of sets is a Segal
$0$-stack, which it is (one can see this by applying \ref{0critere} twice
more). This proves that $Ind(A')$ is a Segal $n$-stack. Then by definition it is
the Segal $n$-stack associated to $Ind(A)$. In the situation of the corollary,
one starts from a Segal $n$-prestack $A$ which is object-by-object
$n$-truncated. It follows that $A$ is equivalent to a prestack $Ind(T)$ where
$T$ is a non-Segal $n$-prestack (and whose values are $n$-categories) and the
preceding discussion proves the corollary.
\eop

{\em Exercise:} If $A$ is a non-Segal $n$-prestack (or a Segal one but
$n$-truncated) then one obtains the associated stack $A'$ from $A$ by applying
the construction of Proposition \ref{exo} $n+2$ times.
